\chapter{Hito 11: Preparación del Proyecto para GitHub y Despliegue}

\section{Objetivo}

El objetivo de este hito es preparar el proyecto para ser publicado en GitHub y posteriormente desplegado en Streamlit Community Cloud.

Hasta este punto ya hemos construido:

\begin{verbatim}
CSV
|
V
Python
|
V
DuckDB
|
V
staging
|
V
core
|
V
mart
|
V
Streamlit
\end{verbatim}

Ahora debemos organizar el repositorio para que cualquier persona pueda entender, instalar y ejecutar el proyecto.

Al finalizar este hito el estudiante será capaz de:

\begin{itemize}
\item Revisar la estructura final del repositorio.
\item Crear un \texttt{README.md} profesional.
\item Actualizar el archivo \texttt{requirements.txt}.
\item Preparar el archivo \texttt{.gitignore}.
\item Validar que la aplicación se ejecuta desde cero.
\item Preparar el proyecto para publicación en GitHub.
\end{itemize}

\section{Resultado esperado}

Al finalizar este hito el proyecto deberá estar listo para subirse a GitHub.

La estructura esperada será:

\begin{verbatim}
dwh_macrofinanciero_rd/

+-- app/
|   -- streamlit_app.py

+-- data/
|   +-- raw/
|   +-- processed/
|   -- warehouse/

+-- docs/
|   -- book/

+-- etl/
|   +-- db.py
|   +-- load_staging.py
|   +-- load_dimensions.py
|   +-- load_facts.py
|   +-- load_marts.py
|   -- run_pipeline.py

+-- sql/
|   +-- 01_create_schemas.sql
|   +-- 02_staging_ddl.sql
|   +-- 03_dimensions_ddl.sql
|   +-- 04_load_dimensions.sql
|   +-- 05_facts_ddl.sql
|   +-- 06_load_facts.sql
|   -- 07_marts.sql

+-- tests/

+-- README.md
+-- requirements.txt
+-- .gitignore
+-- .env.example
\end{verbatim}

\section{Paso 1: Revisar estructura del proyecto}

Desde la raíz del proyecto ejecutar:

\begin{lstlisting}[language=bash]
tree /f
\end{lstlisting}

En caso de no tener disponible el comando \texttt{tree}, usar:

\begin{lstlisting}[language=bash]
dir
\end{lstlisting}

Verificar que las carpetas principales existen:

\begin{itemize}
\item \texttt{app}
\item \texttt{data}
\item \texttt{docs}
\item \texttt{etl}
\item \texttt{sql}
\item \texttt{tests}
\end{itemize}

\section{Paso 2: Revisar el archivo \texttt{.gitignore}}

El archivo \texttt{.gitignore} evita subir al repositorio archivos innecesarios, privados o pesados.

Debe contener:

\begin{lstlisting}
.venv/
.env

**pycache**/
*.pyc
*.pyo

.vscode/

data/processed/
data/warehouse/*.duckdb
data/warehouse/*.wal

.DS_Store
Thumbs.db
\end{lstlisting}

\section{Nota importante sobre DuckDB}

El archivo:

\begin{verbatim}
data/warehouse/macrofinanzas.duckdb
\end{verbatim}

no debe subirse a GitHub si el proyecto crecerá en tamaño.

En su lugar, el usuario debe poder reconstruir la base ejecutando:

\begin{lstlisting}[language=bash]
python etl/run_pipeline.py
\end{lstlisting}

Esta es una práctica recomendable porque permite reproducibilidad.

\section{Paso 3: Revisar \texttt{.env.example}}

El archivo:

\begin{verbatim}
.env.example
\end{verbatim}

debe documentar la ruta esperada de la base DuckDB.

Contenido sugerido:

\begin{lstlisting}
DUCKDB_PATH=data/warehouse/macrofinanzas.duckdb
\end{lstlisting}

El archivo real:

\begin{verbatim}
.env
\end{verbatim}

no debe subirse a GitHub.

\section{Paso 4: Actualizar \texttt{requirements.txt}}

Con el entorno virtual activo:

\begin{lstlisting}[language=bash]
pip freeze > requirements.txt
\end{lstlisting}

Verificar que incluya, al menos:

\begin{itemize}
\item \texttt{duckdb}
\item \texttt{pandas}
\item \texttt{python-dotenv}
\item \texttt{streamlit}
\item \texttt{plotly}
\end{itemize}

\section{Paso 5: Crear un README profesional}

Actualizar el archivo:

\begin{verbatim}
README.md
\end{verbatim}

con el siguiente contenido sugerido:

\begin{lstlisting}

# Data Warehouse Macrofinanciero RD

Proyecto aplicado de ingeniería de datos para construir una plataforma de inteligencia macrofinanciera utilizando Python, DuckDB y Streamlit.

## Objetivo

Construir un flujo completo de datos que permita:

* Ingerir datos desde archivos CSV.
* Almacenarlos en DuckDB.
* Organizar la información en capas staging, core y mart.
* Construir indicadores macrofinancieros.
* Publicar un dashboard interactivo en Streamlit.

## Arquitectura

CSV -> Python -> DuckDB -> Staging -> Core -> Mart -> Streamlit

## Tecnologías

* Python
* DuckDB
* Pandas
* Streamlit
* Plotly
* Git
* GitHub

## Estructura del proyecto

app/
data/
docs/
etl/
sql/
tests/

## Cómo ejecutar el proyecto

1. Crear entorno virtual

python -m venv .venv

2. Activar entorno virtual

.venv\Scripts\activate

3. Instalar dependencias

pip install -r requirements.txt

4. Ejecutar pipeline

python etl/run_pipeline.py

5. Abrir dashboard

streamlit run app/streamlit_app.py

## Autor

Francisco Ramírez
\end{lstlisting}

\section{Paso 6: Validar ejecución completa}

Desde la raíz del proyecto ejecutar:

\begin{lstlisting}[language=bash]
python etl/run_pipeline.py
\end{lstlisting}

Verificar que se ejecutan correctamente:

\begin{itemize}
\item \texttt{load_staging.py}
\item \texttt{load_dimensions.py}
\item \texttt{load_facts.py}
\item \texttt{load_marts.py}
\end{itemize}

\section{Paso 7: Validar dashboard local}

Ejecutar:

\begin{lstlisting}[language=bash]
streamlit run app/streamlit_app.py
\end{lstlisting}

Verificar:

\begin{itemize}
\item La aplicación abre en el navegador.
\item Los KPIs se muestran correctamente.
\item Los gráficos se renderizan.
\item Las tablas se visualizan.
\item No aparecen errores de conexión con DuckDB.
\end{itemize}

\section{Paso 8: Inicializar Git}

Si el repositorio aún no está inicializado:

\begin{lstlisting}[language=bash]
git init
\end{lstlisting}

Verificar estado:

\begin{lstlisting}[language=bash]
git status
\end{lstlisting}

\section{Paso 9: Agregar archivos al repositorio}

Ejecutar:

\begin{lstlisting}[language=bash]
git add .
\end{lstlisting}

Revisar qué archivos serán incluidos:

\begin{lstlisting}[language=bash]
git status
\end{lstlisting}

Confirmar que no aparecen:

\begin{itemize}
\item \texttt{.env}
\item \texttt{.venv/}
\item \texttt{macrofinanzas.duckdb}
\end{itemize}

\section{Paso 10: Crear commit}

Ejecutar:

\begin{lstlisting}[language=bash]
git commit -m "Prepare project for GitHub deployment"
\end{lstlisting}

Verificar historial:

\begin{lstlisting}[language=bash]
git log --oneline
\end{lstlisting}

\section{Paso 11: Crear repositorio en GitHub}

Entrar a GitHub y crear un repositorio nuevo.

Nombre sugerido:

\begin{verbatim}
dwh_macrofinanciero_rd
\end{verbatim}

Seleccionar:

\begin{itemize}
\item Repositorio público o privado.
\item No inicializar con README si ya existe uno local.
\item No agregar \texttt{.gitignore} desde GitHub si ya existe localmente.
\end{itemize}

\section{Paso 12: Conectar repositorio local con GitHub}

GitHub mostrará instrucciones similares a:

\begin{lstlisting}[language=bash]
git remote add origin https://github.com/usuario/dwh_macrofinanciero_rd.git
git branch -M main
git push -u origin main
\end{lstlisting}

Ejecutarlas desde la raíz del proyecto.

\section{Paso 13: Verificar publicación}

Entrar al repositorio en GitHub y verificar que aparecen:

\begin{itemize}
\item \texttt{README.md}
\item \texttt{requirements.txt}
\item carpeta \texttt{app}
\item carpeta \texttt{etl}
\item carpeta \texttt{sql}
\item carpeta \texttt{docs}
\end{itemize}

\section{Paso 14: Preparar archivo para despliegue}

Streamlit Community Cloud utilizará:

\begin{verbatim}
requirements.txt
\end{verbatim}

para instalar dependencias y:

\begin{verbatim}
app/streamlit_app.py
\end{verbatim}

como archivo principal de la aplicación.

Verificar que en GitHub existan ambos archivos.

\section{Paso 15: Consideración importante sobre datos}

Si el archivo \texttt{macrofinanzas.duckdb} no se sube a GitHub, entonces la app desplegada no tendrá base de datos a menos que:

\begin{itemize}
\item La app construya la base al iniciar.
\item El pipeline se ejecute antes del despliegue.
\item Se suba una versión pequeña de la base DuckDB.
\end{itemize}

Para esta primera versión educativa, hay dos opciones.

\subsection{Opción A: Subir una base pequeña}

Permite que la app funcione inmediatamente en Streamlit Cloud.

En este caso se puede permitir temporalmente:

\begin{verbatim}
data/warehouse/macrofinanzas.duckdb
\end{verbatim}

en GitHub, siempre que sea pequeño y no contenga información sensible.

\subsection{Opción B: Reconstruir la base}

Modificar la app para que, si no encuentra la base, ejecute el pipeline automáticamente.

Esta opción es más robusta, pero requiere más cuidado.

\section{Recomendación para este proyecto}

Para la primera publicación, se recomienda la Opción A:

\begin{quote}
Subir una base DuckDB pequeña, construida con datos simulados.
\end{quote}

Más adelante se puede migrar a una arquitectura más robusta donde la base se reconstruya automáticamente.

\section{Buenas prácticas}

Durante este hito seguiremos estas reglas:

\begin{enumerate}
\item No subir credenciales al repositorio.
\item Mantener \texttt{requirements.txt} actualizado.
\item Documentar claramente cómo ejecutar el proyecto.
\item Asegurar que el pipeline reconstruye la base.
\item Usar datos simulados para despliegues públicos.
\end{enumerate}

\section{Checklist de validación}

\begin{itemize}
\item[$\Box$] Archivo \texttt{README.md} actualizado.
\item[$\Box$] Archivo \texttt{requirements.txt} actualizado.
\item[$\Box$] Archivo \texttt{.gitignore} revisado.
\item[$\Box$] Pipeline ejecuta correctamente.
\item[$\Box$] Dashboard abre localmente.
\item[$\Box$] Repositorio Git inicializado.
\item[$\Box$] Primer commit de despliegue creado.
\item[$\Box$] Repositorio publicado en GitHub.
\end{itemize}

\section{Entregable}

Al finalizar este hito el estudiante deberá disponer de:

\begin{itemize}
\item Repositorio GitHub profesional.
\item Proyecto documentado.
\item Dashboard ejecutable localmente.
\item Pipeline reproducible.
\item Proyecto listo para desplegar en Streamlit Cloud.
\end{itemize}

\section{Próximo hito}

En el Hito 12 desplegaremos la aplicación en Streamlit Community Cloud.

Aprenderemos a:

\begin{itemize}
\item Conectar GitHub con Streamlit Cloud.
\item Configurar el archivo principal de la app.
\item Gestionar dependencias.
\item Resolver errores comunes de despliegue.
\item Obtener una URL pública del dashboard.
\end{itemize}
